\section{问题重述}

随着数据中心规模和AI算力需求的快速增长，全球算力中心用电量已占全社会用电量的1.6\%。预计2030年超过全社会用电量的5\%，因此算电协同已于2026年首次纳入国家级新基建工程。实现智能化任务调度与储能协同，便可在时间和空间上动态匹配算力需求与电力供给，而这正是实现数据中心绿色低碳转型的关键。基于50\,000条任务轨迹与2\,407小时逐时电力、碳强度数据，以及对六个典型区域的调度问题建立数学模型，我们发现六区域特征差异明显，其中：东部A、B、C区负荷高，电价与碳强度偏高，时延却较低。西部D区GPU容量最大。E、F区分别富集光伏与风电，电价与碳强度低，新能源装机高。数据时域为0--2406h，其中主时域0--2399h承担任务到达与调度，而2400--2405h是缓冲完成期，2406h仅做纯电力结算。任务按类型分为三类，实时推理、批量推理和AI训练，其中实时推理到达即开工，时延上限20\,ms；批量推理的时延上限放宽到80\,ms；AI训练的时延上限为150\,ms，AI训练是低碳调度的主要调节对象。所有任务均不允许抢占、拆分或中途迁移，必须在时点2406之前完成。本题要求依次求解四个子问题。

\begin{enumerate}
    \item \textbf{GPU需求统计分析与短期预测，以及基础算力调度。} 子问题一包含两个相互独立的小问，可分开求解。第一问分析整体与各区域的GPU逐小时需求特征，并建立短期预测模型，预测第2376--2399小时的逐时GPU需求，这24小时构成评估时段。第二问建立基础算力调度模型，在GPU容量、时延与完成时限的约束下确定各任务的执行区域与开工时段，使GPU利用率保持均衡，且全部任务按时完成。

    \item \textbf{碳感知任务调度模型。} 在子问题一基础上引入跨区域迁移与开工时段自由度，并综合考虑电价、碳强度、PUE、时延与可用新能源，从而建立以运行成本与碳排放为核心指标的多目标调度模型。其中，实时推理的时延上限为20\,ms，只能在A/B/C之间或E/F之间互迁。批量推理的时延上限为80\,ms，A至F的单向路径被排除。而AI训练不受区域限制，可在六区域间自由迁移。本题要求给出最优任务分配方案，同时量化该迁移方案相对于零迁移基线，在成本上的变化和在碳减排上的收益。

    \item \textbf{储能协同优化模型。} 固定子问题二的IT负荷，在六区域独立配置储能的条件下建立充放电、购售电与新能源分配的协同优化模型。储能须满足以下约束：容量上下限、同一时段不可同时充放电、SOC递推等，其中初始SOC为45\%容量，充电效率$\eta_c\approx0.93$，放电效率$\eta_d\approx0.92$。本题要求给出各区域储能最优运行策略、逐时购售电方案与新能源消纳结果。

    \item \textbf{多区域算--储--电协同优化与场景分析。} 子问题四统一优化任务调度与储能运行，在六区域联合框架下构建多目标协同模型。该模型放开任务调度与储能运行的全部自由，以运行成本、碳排放、网络时延、QoS、新能源利用率与区域峰值净购电功率六项指标为评价维度。我们针对不同电价、碳约束强度与新能源可用量场景，比较出最优策略，同时分析关键参数灵敏度。
\end{enumerate}
